On compare la diffusion Compton, à un angle de \qty{\v{th}}{\degree}, de deux types de photons sur des électrons libres : de la lumière visible de \qty{\v{lv}}{nm} et des rayons X de \qty{\v{lx}}{pm}.

\begin{enumerate}[label=\alph*)]
    \item Déterminez le décalage de longueur d'onde $\Delta\lambda$. Dépend-il de la longueur d'onde incidente ?
    \item Déterminez la variation relative $\Delta\lambda/\lambda$ pour la lumière visible.
    \item Déterminez la variation relative $\Delta\lambda/\lambda$ pour les rayons X. Expliquez pourquoi l'effet Compton a été découvert avec des rayons X et non avec de la lumière visible.
\end{enumerate}
